\begin{thebibliography}{99}
\setlength{\itemsep}{2pt}

\bibitem[Afrifa-Yamoah et~al.(2026)]{afrifa2026}
Afrifa-Yamoah, E., Awuah-Mensah, Y. K., Bradshaw, S. and Mueller, U. (2026). Climate change impacts on fish-related ecosystem services: synthesis, network analysis, and strategic pathways for coordinated adaptation. \textit{Reviews in Fish Biology and Fisheries}, 36, 51.

\bibitem[Arregu\'in-S\'anchez(1996)]{arreguin1996}
Arregu\'in-S\'anchez, F. (1996). Catchability: a key parameter for fish stock assessment. \textit{Reviews in Fish Biology and Fisheries}, 6(2), 221--242.

\bibitem[Awuah-Mensah and Afrifa-Yamoah(2026)]{awuah2026}
Awuah-Mensah, Y. K. and Afrifa-Yamoah, E. (2026). Code: Declined, displaced, or just harder to catch? A simulation framework for attributing marine heatwave impacts on fished populations from catch per unit effort (Version 2.0.0) [Computer software]. \textit{Zenodo}. \url{https://doi.org/10.5281/zenodo.23629610}

\bibitem[Babcock et~al.(2019)]{babcock2019}
Babcock, R. C., Bustamante, R. H., Fulton, E. A. et al. (2019). Severe continental-scale impacts of climate change are happening now: extreme climate events impact marine habitat forming communities along 45\% of Australia's coast. \textit{Frontiers in Marine Science}, 6, 411.

\bibitem[Caputi et~al.(2016)]{caputi2016}
Caputi, N., Kangas, M., Denham, A., Feng, M., Pearce, A., Hetzel, Y. and Chandrapavan, A. (2016). Management adaptation of invertebrate fisheries to an extreme marine heat wave event at a global warming hot spot. \textit{Ecology and Evolution}, 6(11), 3583--3593.

\bibitem[Caputi et~al.(2019)]{caputi2019}
Caputi, N., Kangas, M., Chandrapavan, A., Hart, A., Feng, M., Marin, M. and de Lestang, S. (2019). Factors affecting the recovery of invertebrate stocks from the 2011 Western Australian extreme marine heatwave. \textit{Frontiers in Marine Science}, 6, 484.

\bibitem[Chaikin et~al.(2024)]{chaikin2024}
Chaikin, S., Riva, F., Marshall, K. E., Lessard, J.-P. and Belmaker, J. (2024). Marine fishes experiencing high-velocity range shifts may not be climate change winners. \textit{Nature Ecology \& Evolution}, 8(5), 936--946.

\bibitem[Cheung and Fr\"olicher(2020)]{cheung2020}
Cheung, W. W. L. and Fr\"olicher, T. L. (2020). Marine heatwaves exacerbate climate change impacts for fisheries in the northeast Pacific. \textit{Scientific Reports}, 10, 6678.

\bibitem[Craven and Wahba(1979)]{craven1979}
Craven, P. and Wahba, G. (1979). Smoothing noisy data with spline functions: estimating the correct degree of smoothing by the method of generalized cross-validation. \textit{Numerische Mathematik}, 31(4), 377--403.

\bibitem[Ducharme-Barth et~al.(2022)]{ducharme2022}
Ducharme-Barth, N. D., Gr\"uss, A., Vincent, M. T., Kiyofuji, H., Aoki, Y., Pilling, G., Hampton, J. and Thorson, J. T. (2022). Impacts of fisheries-dependent spatial sampling patterns on catch-per-unit-effort standardization: a simulation study and fishery application. \textit{Fisheries Research}, 246, 106169.

\bibitem[Erisman et~al.(2011)]{erisman2011}
Erisman, B. E., Allen, L. G., Claisse, J. T., Pondella, D. J., Miller, E. F. and Murray, J. H. (2011). The illusion of plenty: hyperstability masks collapses in two recreational fisheries that target fish spawning aggregations. \textit{Canadian Journal of Fisheries and Aquatic Sciences}, 68(10), 1705--1716.

\bibitem[Feng et~al.(2013)]{feng2013}
Feng, M., McPhaden, M. J., Xie, S.-P. and Hafner, J. (2013). La Ni\~{n}a forces unprecedented Leeuwin Current warming in 2011. \textit{Scientific Reports}, 3, 1277.

\bibitem[Fredston et~al.(2023)]{fredston2023}
Fredston, A. L., Cheung, W. W. L., Fr\"olicher, T. L. et al. (2023). Marine heatwaves are not a dominant driver of change in demersal fishes. \textit{Nature}, 621(7978), 324--329.

\bibitem[Fretwell and Lucas(1969)]{fretwell1969}
Fretwell, S. D. and Lucas, H. L. (1969). On territorial behavior and other factors influencing habitat distribution in birds. \textit{Acta Biotheoretica}, 19(1), 16--36.

\bibitem[Fr\"olicher et~al.(2018)]{frolicher2018}
Fr\"olicher, T. L., Fischer, E. M. and Gruber, N. (2018). Marine heatwaves under global warming. \textit{Nature}, 560(7718), 360--364.

\bibitem[Fuster-Alonso et~al.(2024)]{fuster2024}
Fuster-Alonso, A., Conesa, D., Pennino, M. G. et al. (2024). Spatio-temporal sampling designs and dependence: implications for abundance indices. \textit{ICES Journal of Marine Science}, 81(5), 984--995.

\bibitem[Harley et~al.(2001)]{harley2001}
Harley, S. J., Myers, R. A. and Dunn, A. (2001). Is catch-per-unit-effort proportional to abundance? \textit{Canadian Journal of Fisheries and Aquatic Sciences}, 58(9), 1760--1772.

\bibitem[Hobday et~al.(2016)]{hobday2016}
Hobday, A. J., Alexander, L. V., Perkins, S. E. et al. (2016). A hierarchical approach to defining marine heatwaves. \textit{Progress in Oceanography}, 141, 227--238.

\bibitem[Horvitz and Thompson(1952)]{horvitz1952}
Horvitz, D. G. and Thompson, D. J. (1952). A generalization of sampling without replacement from a finite universe. \textit{Journal of the American Statistical Association}, 47(260), 663--685.

\bibitem[Hoyle et~al.(2024)]{hoyle2024}
Hoyle, S. D., Campbell, R. A., Ducharme-Barth, N. D. et al. (2024). Catch per unit effort modelling for stock assessment: a summary of good practices. \textit{Fisheries Research}, 269, 106860.

\bibitem[Karp et~al.(2023)]{karp2023}
Karp, M. A., Brodie, S., Smith, J. A. et al. (2023). Projecting species distributions using fishery-dependent data. \textit{Fish and Fisheries}, 24(1), 71--92.

\bibitem[Maunder and Punt(2004)]{maunder2004}
Maunder, M. N. and Punt, A. E. (2004). Standardizing catch and effort data: a review of recent approaches. \textit{Fisheries Research}, 70(2--3), 141--159.

\bibitem[Maunder et~al.(2006)]{maunder2006}
Maunder, M. N., Sibert, J. R., Fonteneau, A., Hampton, J., Kleiber, P. and Harley, S. J. (2006). Interpreting catch per unit effort data to assess the status of individual stocks and communities. \textit{ICES Journal of Marine Science}, 63(8), 1373--1385.

\bibitem[Maunder et~al.(2020)]{maunder2020}
Maunder, M. N., Thorson, J. T., Xu, H. et al. (2020). The need for spatio-temporal modeling to determine catch-per-unit effort based indices of abundance and associated composition data for inclusion in stock assessment models. \textit{Fisheries Research}, 229, 105594.

\bibitem[Newey and West(1987)]{newey1987}
Newey, W. K. and West, K. D. (1987). A simple, positive semi-definite, heteroskedasticity and autocorrelation consistent covariance matrix. \textit{Econometrica}, 55(3), 703--708.

\bibitem[Neyman and Scott(1948)]{neyman1948}
Neyman, J. and Scott, E. L. (1948). Consistent estimates based on partially consistent observations. \textit{Econometrica}, 16(1), 1--32.

\bibitem[Oliver et~al.(2018)]{oliver2018}
Oliver, E. C. J., Donat, M. G., Burrows, M. T. et al. (2018). Longer and more frequent marine heatwaves over the past century. \textit{Nature Communications}, 9, 1324.

\bibitem[Pearce and Feng(2013)]{pearce2013}
Pearce, A. F. and Feng, M. (2013). The rise and fall of the ``marine heat wave'' off Western Australia during the summer of 2010/2011. \textit{Journal of Marine Systems}, 111--112, 139--156.

\bibitem[S\"arndal et~al.(1992)]{sarndal1992}
S\"arndal, C.-E., Swensson, B. and Wretman, J. (1992). \textit{Model Assisted Survey Sampling}. Springer.

\bibitem[Smale et~al.(2019)]{smale2019}
Smale, D. A., Wernberg, T., Oliver, E. C. J. et al. (2019). Marine heatwaves threaten global biodiversity and the provision of ecosystem services. \textit{Nature Climate Change}, 9(4), 306--312.

\bibitem[Thorson(2019)]{thorson2019}
Thorson, J. T. (2019). Guidance for decisions using the Vector Autoregressive Spatio-Temporal (VAST) package in stock, ecosystem, habitat and climate assessments. \textit{Fisheries Research}, 210, 143--161.

\bibitem[Thorson and Barnett(2017)]{thorson2017}
Thorson, J. T. and Barnett, L. A. K. (2017). Comparing estimates of abundance trends and distribution shifts using single- and multispecies models of fishes and biogenic habitat. \textit{ICES Journal of Marine Science}, 74(5), 1311--1321.

\bibitem[Thorson et~al.(2020)]{thorson2020}
Thorson, J. T., Maunder, M. N. and Punt, A. E. (2020). The development of spatio-temporal models of fishery catch-per-unit-effort data to derive indices of relative abundance. \textit{Fisheries Research}, 230, 105611.

\bibitem[Thorson et~al.(2016a)]{thorson2016a}
Thorson, J. T., Pinsky, M. L. and Ward, E. J. (2016a). Model-based inference for estimating shifts in species distribution, area occupied and centre of gravity. \textit{Methods in Ecology and Evolution}, 7(8), 990--1002.

\bibitem[Thorson et~al.(2016b)]{thorson2016b}
Thorson, J. T., Rindorf, A., Gao, J., Hanselman, D. H. and Winker, H. (2016b). Density-dependent changes in effective area occupied for sea-bottom-associated marine fishes. \textit{Proceedings of the Royal Society B: Biological Sciences}, 283(1840), 20161853.

\bibitem[Thorson et~al.(2015)]{thorson2015}
Thorson, J. T., Shelton, A. O., Ward, E. J. and Skaug, H. J. (2015). Geostatistical delta-generalized linear mixed models improve precision for estimated abundance indices for West Coast groundfishes. \textit{ICES Journal of Marine Science}, 72(5), 1297--1310.

\bibitem[Welch et~al.(2023)]{welch2023}
Welch, H., Savoca, M. S., Brodie, S. et al. (2023). Impacts of marine heatwaves on top predator distributions are variable but predictable. \textit{Nature Communications}, 14, 5188.

\bibitem[Wilberg et~al.(2010)]{wilberg2010}
Wilberg, M. J., Thorson, J. T., Linton, B. C. and Berkson, J. (2010). Incorporating time-varying catchability into population dynamic stock assessment models. \textit{Reviews in Fisheries Science}, 18(1), 7--24.

\bibitem[Wernberg et~al.(2016)]{wernberg2016}
Wernberg, T., Bennett, S., Babcock, R. C. et al. (2016). Climate-driven regime shift of a temperate marine ecosystem. \textit{Science}, 353(6295), 169--172.

\bibitem[Wernberg et~al.(2013)]{wernberg2013}
Wernberg, T., Smale, D. A., Tuya, F., Thomsen, M. S., Langlois, T. J., de Bettignies, T., Bennett, S. and Rousseaux, C. S. (2013). An extreme climatic event alters ecosystem structure in a global biodiversity hotspot. \textit{Nature Climate Change}, 3(1), 78--82.

\end{thebibliography}